\makeabstract
    {
    Borderline Personality Disorder and Interpersonal Emotion Regulation Strategies
    }
    {
    Ani Bailin\textsuperscript{1}, Mabel Shanahan\textsuperscript{1}, Elizabeth T. Kneeland\textsuperscript{1}
    }
    {
    \textsuperscript{1} Department of Psychology, Amherst College, Amherst, MA
    }
    {
    This study investigates the associations between BPD symptoms and use of interpersonal emotion regulation strategies. The study examines correlations between PAI-BOR scores and use of interpersonal rumination, expressive suppression, acceptance, and cognitive reappraisal. 
    }
    {
    This study analyzed previous data from a larger study on emotion regulation and emotion beliefs. All data was self-reported and collected via Qualtrics survey. Participants (N=111; mean age=28; 73.9\% female; 55.9\% white) in this sample completed the Personality Assessment Inventory-Borderline (PAI-BOR) and the Emotion Regulation Strategy Scale (ERSS). 
    }
    {
    Correlation analysis reveals that participants with greater BPD symptom severity were more likely to engage in interpersonal expressive suppression. Participants who scored higher in BPD affective instability were less likely to engage in interpersonal acceptance. Interestingly, there was no correlation between BPD symptoms and use of interpersonal cognitive reappraisal. T-test analysis reveals that participants who met BPD diagnostic criteria were significantly less likely to use interpersonal rumination.
    }
    {
    This study concludes that greater BPD symptom severity is linked to greater use of interpersonal expressive suppression and lesser use of interpersonal acceptance. Those who met diagnostic criteria for BPD were less likely to ruminate with other people. These findings improve our understanding of BPD and provide insights toward developing more effective strategies in managing BPD symptoms. Future research might include a larger sample size and a clinical BPD sample. 
    }

    \newpage
    